% !TEX program = lualatex
\documentclass[11pt,a4paper]{ltjsarticle}
\usepackage{amsmath,amssymb,amsthm}
\usepackage{geometry}
\geometry{margin=1in}
\usepackage{enumitem}
\usepackage{indentfirst}

%-------------------------------------------------
%  マクロ定義
%-------------------------------------------------
\newcommand{\Mdelta}{M_\delta}

%-------------------------------------------------
%  定理環境
%-------------------------------------------------
\theoremstyle{definition}
\newtheorem{definition}{定義}[section]
\newtheorem{proposition}{命題}[section]
\newtheorem{corollary}{帰結}[section]
\newtheorem{remark}{注記}[section]

\begin{document}

\title{搾取モナド：\\
圏論的フレームワークによる富の集約・国家・\\
二元論社会の構造的分析}
\author{千葉将臣}
\date{2026-05-24}
\maketitle

\begin{abstract}
本稿は、モナドおよび随伴関手の圏論的構造を用いて「搾取モナド」という
分析的フレームワークを定式化し、富の一極集中・国家と戦争の発生・
宗教的構造との関係を記述する。
本フレームワークは圏論的構造をアナロジーとして用いた構造的記述の道具として
位置づけられるものであり、厳密な数学的定理の体系ではない。
個々の道徳的評価や意図に依存せず、構造それ自体から搾取の発生条件・持続条件・
限界を記述することが本稿の目的である。
なお本稿において「搾取」はいかなる価値判断も含まない構造的概念として用いる。
\end{abstract}

\tableofcontents
\newpage

%=============================================================
\section{搾取モナドの定式化}
%=============================================================

\subsection{基礎概念の定義}

\begin{definition}[截断]
対象 $A$ の多層的・連続的・フラクタル的な差異を、
操作可能な自己同一性 $\mathrm{id}_A$ へ圧縮する操作を\textbf{截断}と呼ぶ。
特に、$A$ と $\neg A$ の二値化によって中間状態を切り捨てる截断を
\textbf{二元論的截断}と呼ぶ。
\end{definition}

\begin{definition}[粒度差 $\delta$]
主体 $X$ と $Y$ が同一の価値対象 $A$ にアクセスするとき、
両者のアクセス可能性・可視性・操作可能性・制度的承認の差を
$\delta(X, Y; A)$ と書く。文脈が明らかな場合これを単に $\delta$ と書く。
$\delta = 0$ は等アクセス状態、$\delta > 0$ は非対称なアクセス構造を表す。
\end{definition}

\begin{definition}[搾取文脈 $\Mdelta$]
粒度差 $\delta$ を持つ文脈を $\Mdelta$ と表記する。
$\Mdelta A$ は「$A$ が制度・法・信用・貨幣・プラットフォーム等の文脈を
経由してアクセスされる状態」を表す。
\end{definition}

\begin{definition}[搾取モナド]
搾取モナドは三つ組 $(\Mdelta,\, \eta,\, \mu)$ として定義される。
\begin{align}
\eta_A &: A \to \Mdelta A \\
\mu_A &: \Mdelta\Mdelta A \to \Mdelta A
\end{align}
$\eta_A$ は価値や主体を $\Mdelta$ 文脈に編入する操作であり、
$\mu_A$ は多重化した文脈を一つの制度的文脈へ統合する操作である。
この三つ組が以下のモナド則を満たすとき搾取モナドと呼ぶ。
\begin{align}
\mu \circ M_\delta\mu &= \mu \circ \mu_{M_\delta} \\
\mu \circ M_\delta\eta &= \mu \circ \eta_{M_\delta} = \mathrm{id}_{M_\delta}
\end{align}
\end{definition}

\subsection{モナド則の社会的解釈}

\begin{remark}
モナド則は搾取の構造的不可避性を示す。
結合則は $\delta$ の累積が括り方に依存せず一定であることを保証し、
単位元則は操作によらず $\delta$ が保存されることを示す。
搾取モナドの各操作は局所的に見れば正当な手続きであり、
モナド則自体はいかなる倫理的問題も含まない。
搾取構造は操作の内容ではなく、\textbf{文脈 $\Mdelta$ の構造として埋め込まれる}。
\end{remark}

\subsection{搾取構造の成立条件}

\begin{proposition}[搾取構造の成立条件]
中立な搾取モナドが搾取構造として作動するのは以下の三条件が重なったときである。
\begin{enumerate}
\item \textbf{$\delta$ の非対称な固定}：
  粒度差の設定権限が $\Mdelta$ の受益者に集中する。
\item \textbf{截断の不可視化}：
  $\delta \neq 0$ という初期条件への問いが自明視によって封鎖される。
\item \textbf{短絡経路の封鎖}：
  $\Mdelta$ を経由しない $A \to B$ が制度的・技術的・法的に遮断される。
\end{enumerate}
三条件は倫理的意図とは独立に作動する。
\end{proposition}

%=============================================================
\section{搾取モナドと圏の構造}
%=============================================================

\subsection{随伴関手による基礎付け}

\begin{definition}[搾取随伴]
圏 $\mathcal{C}$（粒度差のない基底圏）と圏 $\mathcal{D}$（$\Mdelta$ 文脈を持つ制度圏）の間に
随伴
\[
F : \mathcal{C} \rightleftharpoons \mathcal{D} : G, \quad F \dashv G
\]
を設定し、$\Mdelta = G \circ F$ と置く。
$F$ は価値を制度文脈へ埋め込む関手、$G$ は制度文脈から価値を取り出す関手である。
\end{definition}

\begin{proposition}[単位・余単位の搾取的解釈]
随伴の単位 $\eta: \mathrm{id}_{\mathcal{C}} \Rightarrow GF$ は
価値・主体が $\Mdelta$ 文脈に編入される入口に対応する。
余単位 $\varepsilon: FG \Rightarrow \mathrm{id}_{\mathcal{D}}$ は
$\Mdelta$ 文脈の往復において $\delta$ が徴収・吸収される局面に対応する。
搾取の実質的発生点は余単位 $\varepsilon$ に局所化される。
\end{proposition}

\begin{corollary}[三角恒等式による不可避性]
随伴の三角恒等式
$G\varepsilon \circ \eta_G = \mathrm{id}_G$,
$\varepsilon_F \circ F\eta = \mathrm{id}_F$
は、$\delta$ の徴収がいかなる経路を通っても消去できないことの
圏論的保証を与える。
\end{corollary}

\subsection{クライスリ圏とアイレンベルク・ムーア圏}

\begin{definition}[クライスリ圏 $\mathcal{C}_{\Mdelta}$]
射を $A \to \Mdelta B$ とするクライスリ圏は、
$\delta$ を経由する可能性のある操作の全体、すなわち\textbf{潜在的搾取の圏}である。
二つの射の合成は
\[
g \star f = \mu_C \circ \Mdelta g \circ f
\]
で与えられる。
\end{definition}

\begin{definition}[$\Mdelta$ 代数（アイレンベルク・ムーア圏）]
$\Mdelta$ 代数 $(A, \alpha: \Mdelta A \to A)$ の圏
$\mathcal{C}^{\Mdelta}$ は、$\delta$ の吸収を制度として内面化した主体の圏、
すなわち\textbf{制度化された搾取の圏}である。
国家・金融機関・巨大プラットフォームはこの $\Mdelta$ 代数の
具体的実装として読める。
\end{definition}

\begin{definition}[短絡経路]
文脈 $\Mdelta$ を経由せずに $A$ から $B$ へ到達する射 $h: A \to B$ を
\textbf{短絡経路}と呼ぶ。
DIY・地産地消・直接交換・共同体内互助・局所 AI による知識獲得は
社会的な短絡経路の例である。
\end{definition}

\begin{proposition}[比較関手と制度化]
クライスリ圏からアイレンベルク・ムーア圏への比較関手
$K: \mathcal{C}_{\Mdelta} \to \mathcal{C}^{\Mdelta}$
は、潜在的搾取が制度的搾取へ移行する経路を記述する。
バー・ベック定理の条件が満たされるとき、
社会構造が $\Mdelta$ 代数と圏同値になることは搾取の完全制度化に対応する。
この比較関手が実質的に不可逆になるとき、
搾取は偶発的関係ではなく制度的秩序となる。
\end{proposition}

\begin{remark}
本節の命題は圏論的アナロジーによる構造記述であり、
社会現象の完全な記述を主張するものではない。
$\delta$ の操作的定義と測定可能な代理変数の設定は今後の課題として残る。
\end{remark}

%=============================================================
\section{富の集約の搾取モナドによる説明}
%=============================================================

\subsection{余剰の発生と搾取モナドの起動}

\begin{proposition}[搾取モナドの起動条件]
自給的均衡（生産 $\approx$ 消費）においては $\delta \approx 0$ が
自然条件として成立し、$\Mdelta$ 文脈が生成される素材がないため
搾取モナドは起動しない。
農業・牧畜による余剰（生産 $>$ 消費）の発生は差分 $\delta > 0$ を生み、
余剰の管理権限が $\delta$ の設定権限へ写像されることで搾取モナドが起動する。
\end{proposition}

\begin{proposition}[持続可能性の逆説]
拡大できない構造（自給的均衡）は $\delta \approx 0$ のため搾取モナドが起動せず持続する。
拡大できる構造（農業的余剰）は $\Mdelta$ を起動させて拡大し、
最終的に自然の再生能力を超えて持続しなくなる。
自然への搾取と人間への搾取は同一の余剰発生から同時に起動した構造を持つ。
\end{proposition}

\subsection{一極集中と既得権益層の随伴強化}

\begin{proposition}[一極集中の構造]
$\Mdelta$ 代数 $(A, \alpha)$ を保持する主体は以下の三つの自己維持操作を構造的に実行する。
\begin{enumerate}
\item $\eta$ の普遍性の確保（短絡経路の封鎖）
\item $\varepsilon$ の安定化（貨幣・法制度の維持）
\item 比較関手 $K$ の不可逆化（参入障壁・規制）
\end{enumerate}
これらは意図的な操作である以前に、$\Mdelta$ 代数の自己維持として
構造的に発生する。
\end{proposition}

\subsection{緩和機構としての短絡経路}

\begin{proposition}[DIY・地産地消の緩和効果]
短絡経路の確保は $\delta$ の連鎖距離を縮め $\delta \approx 0$ の
局所的実現として機能する。
これは搾取モナドの緩和であって、$\Mdelta$ の生成条件そのものの解体ではない。
局所 AI・エッジ AI は知識獲得の短絡経路を生成する道具として機能しうるが、
局所 AI 自体が新たな $\Mdelta$ 代数になるリスクも内包する。
\end{proposition}

%=============================================================
\section{国家と搾取モナド}
%=============================================================

\begin{definition}[国家の圏論的記述]
国家とは $\Mdelta$ 代数の最大実装であり、
$\delta$ の設定（税率決定）・徴収（課税）・運用（財政支出）という
三権限を単一主体に集約した $\alpha_{\text{国家}}: \Mdelta A \to A$ として読める。
ウェーバーの「暴力の独占」は $\varepsilon$ の徴収権の単一化として再記述できる。
\end{definition}

\begin{proposition}[戦争の発生条件]
複数の $\Mdelta$ 代数が同一の基底圏に対して異なる随伴構造を主張するとき、
すなわち $\eta_1: A \to M_{\delta_1}A$ と $\eta_2: A \to M_{\delta_2}A$ が
競合するとき、随伴の普遍性への圧力から一方の消滅による
単一随伴の確立へ収束する過程が生じる。これが戦争の構造的論理に対応する。
\end{proposition}

\begin{corollary}[帝国主義と現代プラットフォーム競合]
帝国主義は $M_{\delta_1}$ 代数が $M_{\delta_2}$ 代数を
自身のアイレンベルク・ムーア圏に包摂する操作として記述できる。
現代では国家とプラットフォームが同一の市民・ユーザーに対して、
課税とデータ収集という異なる $\eta$ を競合させている。
\end{corollary}

\begin{corollary}[貨幣と税の一体化]
貨幣による交換が課税と一体化することは、
全ての価値移動が $\Mdelta$ を強制経由する構造であり、
搾取モナドの三条件を一括実装したシステムとなる。
DIY・贈与・物々交換という短絡経路は課税困難性・法的曖昧性により
周縁化される。
\end{corollary}

%=============================================================
\section{二元論社会に必然として発生する搾取モナド}
%=============================================================

\begin{proposition}[二元論の截断機能]
$\mathrm{id}: A \to A$ を操作として機能させるには
対象の多層的・連続的・フラクタル的な差異を $A$ か $\neg A$ かに二値化する截断が前提となる。
二元論はこの截断の最小形態であり、搾取モナドの生成に必要な構造的土台となる。
\end{proposition}

\begin{proposition}[二元論化と搾取モナドの自己強化]
社会の二元論化が進むほど截断が深化し、$\delta$ の固定が容易になることで
搾取モナドの起動条件が整備される。
逆に搾取モナドの強化は截断をさらに推進するという自己強化ループを構成する。
截断は一度発生すると自己維持するため、社会の自発的変化は
搾取モナドで説明できる方向（截断の深化）にのみ進む傾向を持つ。
\end{proposition}

\subsection{マルクスの分析の位置づけ}

\begin{proposition}[マルクスの診断と構造的限界]
マルクスは $\Mdelta$ の作動（剰余価値の抽出）を正確に観察したが、
$\Mdelta$ の生成条件（截断・$\mathrm{id}$ の固定という前提構造）には届かなかった。
処方（生産手段の共有化）は $\Mdelta$ の担い手の交替操作であり、
より大きな $\delta'$ を持つ前衛党という新たな $\Mdelta$ 代数を構造的に生成した
（$M_\delta \to M_{\delta'}$, $\delta' \gg \delta$）。
テキストの正典化によって批判の道具自体が新たな $\Mdelta$ 代数へ転化したことも
失敗を加速した。
\end{proposition}

%=============================================================
\section{一神教・多神教・仏教と搾取モナド}
%=============================================================

\begin{remark}[本節の位置づけと限定]
本節は宗教的教義の評価ではなく、
各宗教体系の構造的特徴と搾取モナドの生成条件との
\textbf{アナロジー的対応}を記述するものである。
以下の分析は宗教の多様な内部差異・神学的伝統・実践の豊かさを捨象した
粗い構造的読みであり、各宗教の歴史的・社会的実態の記述でも
価値判断でもない。
これらは「構造的に何が起きやすいか」という傾向の記述に留まる。
\end{remark}

\subsection{一神教の構造的傾向}

\begin{proposition}[一神教と搾取モナドの構造的親和性]
一神教において唯一神という絶対的 $\mathrm{id}$ の固定は
截断の形而上学的根拠を提供しやすい。
神／人間・聖／俗という二元論的截断は搾取モナドの生成条件と
構造的に整合し、神の代理人・教会・制度という執行主体が
$\Mdelta$ 代数を生成しやすい条件を持つ。
\end{proposition}

\subsection{多神教の構造的傾向}

\begin{proposition}[多神教の緩衝機能]
多神教においては複数の神が $\mathrm{id}$ を分散させ、
截断の絶対化が抑制されやすい。
神々の矛盾・競合・変容が多層的・連続的・フラクタル的な性質の保持として機能し、
搾取モナドの生成条件を一点に集約させない緩衝として働く傾向を持つ。
ただし神々の階層化・主神の特権化・国家祭祀を通じて
$\Mdelta$ 代数が再生成されうる。
\end{proposition}

\subsection{仏教の構造的位置}

\begin{proposition}[仏教と搾取モナドの非接触]
仏教は諸法無我・空という原理によって $\mathrm{id}$ の固定が幻想であることを
体系の出発点とする。これは搾取モナドの第一条件（$\mathrm{id}$ の固定）への介入であり、
搾取モナドが立ち上がる以前の層に位置する。
ただし仏教は個人の $\mathrm{id}$ 固定を解除しうるが
社会的文脈 $\Mdelta$ の設計に直接介入しないため、
搾取モナドからの\textbf{離脱}を可能にするが\textbf{解体}の道具ではない。
\end{proposition}

\begin{corollary}[三者と本フレームワークの対比]
\begin{align}
\text{一神教} &:\ \mathrm{id}\text{ を絶対固定}
  \to \text{截断を聖化}
  \to \Mdelta\text{ を構造的に生成しやすい}\\
\text{多神教} &:\ \mathrm{id}\text{ を分散}
  \to \Mdelta\text{ 生成を抑制、ただし制度的に再生成されうる}\\
\text{仏教}   &:\ \mathrm{id}\text{ 固定を解除}
  \to \Mdelta\text{ の生成条件から降りる}\\
\text{本稿}   &:\ \text{截断を記述}
  \to \Mdelta\text{ の生成構造を可視化する}
\end{align}
\end{corollary}

%=============================================================
\section{結語}
%=============================================================

搾取モナドフレームワークは以下を達成する。

\begin{enumerate}
\item 搾取を意図・倫理・イデオロギーから独立した
  \textbf{構造的記述}として定式化する。
\item 富の集約・国家・戦争・宗教的構造との関係を
  同一フレームワークで記述する。
\item 倫理批判が構造に届かない理由を $\Mdelta$ 代数の自己維持として説明する。
\item 社会主義の失敗と革命の弱体化が同型の構造的矛盾から発生することを記述する。
\end{enumerate}

現状の限界として以下を明記する。

\begin{itemize}
\item $\delta$ の操作的定義と測定可能な代理変数の未整備
\item 随伴アナロジーの形式化（$\delta$ の時間的変化の記述、高次圏への拡張）
\item 宗教・社会現象の記述における内部差異の捨象
\end{itemize}

搾取構造への有効な介入は $\delta$ の設定権限を分散し、截断を可視化し、
短絡経路を回復することにある。倫理批判という不適切な道具に代わる
構造設計の問いとして本フレームワークを展開することが今後の方向である。

\begin{thebibliography}{9}
\bibitem{MacLane1971}
Saunders Mac Lane.
\newblock \emph{Categories for the Working Mathematician}.
\newblock Springer, 1971.

\bibitem{Moggi1991}
Eugenio Moggi.
\newblock Notions of computation and monads.
\newblock \emph{Information and Computation}, 93(1):55--92, 1991.

\bibitem{Hewitt2008}
Carl Hewitt.
\newblock Formalizing common sense for scalable inconsistency-robust
information integration using {Direct Logic} reasoning and the {Actor Model}.
\newblock \emph{arXiv:0812.4852}, 2008.
\end{thebibliography}

\end{document}
